% V tomto souboru se nastavují téměř veškeré informace, proměnné mezi studenty:
% jméno, název práce, pohlaví atd.
% Tento soubor je SDÍLENÝ mezi textem práce a prezentací k obhajobě -- netřeba něco nastavovat na dvou místech.

\usepackage[
%%% Z následujících voleb jazyka lze použít pouze jednu
  czech-english,		% originální jazyk je čeština, překlad je anglicky (výchozí)
  %english-czech,	% originální jazyk je angličtina, překlad je česky
  %slovak-english,	% originální jazyk je slovenština, překlad je anglicky
  %english-slovak,	% originální jazyk je angličtina, překlad je slovensky
%
%%% Z následujících voleb typu práce lze použít pouze jednu
  %semestral,		  % semestrální práce (výchozí)
  %bachelor,			%	bakalářská práce
  master,			  % diplomová práce
  %treatise,			% pojednání o disertační práci
  %doctoral,			% disertační práce
%
%%% Z následujících voleb zarovnání objektů lze použít pouze jednu
%  left,				  % rovnice a popisky plovoucích objektů budou zarovnány vlevo
	center,			    % rovnice a popisky plovoucích objektů budou zarovnány na střed (vychozi)
%
%%% Níže uvedený přepinač 'electronic' lze použít pro generování elektronické verze práce; pokud je aktivní, vnější a vnitřní okraj sazebního obrazce budou shodné pro liché i sudé stránky.
% Pozor, neplést si s volbou oneside/twoside!
% Pozor, pro tiskovou verzi nechejte vypnuté!
%	electronic			
]{thesis}   % Balíček pro sazbu studentských prací


%%% Jméno a příjmení autora ve tvaru
%  [tituly před jménem]{Křestní}{Příjmení}[tituly za jménem]
% Pokud osoba nemá titul před/za jménem, smažte celý řetězec '[...]'
\author[Bc.]{Dušan}{Horák}

%%% Identifikační číslo autora (VUT ID)
\butid{230076}

%%% Pohlaví autora/autorky
% (nepoužije se ve variantě english-czech ani english-slovak)
% Číselná hodnota: 1...žena, 0...muž
\gender{0}

%%% Jméno a příjmení vedoucího/školitele včetně titulů
%  [tituly před jménem]{Křestní}{Příjmení}[tituly za jménem]
% Pokud osoba nemá titul před/za jménem, smažte celý řetězec '[...]'
\advisor[Ing.]{Libor}{Veselý}[Ph.D.]

%%% Jméno a příjmení oponenta včetně titulů
%  [tituly před jménem]{Křestní}{Příjmení}[tituly za jménem]
% Pokud osoba nemá titul před/za jménem, smažte celý řetězec '[...]'
% Nastavení oponenta se uplatní pouze v prezentaci k obhajobě;
% v případě, že nechcete, aby se na titulním snímku prezentace zobrazoval oponent, pouze příkaz zakomentujte;
% u obhajoby semestrální práce se oponent nezobrazuje (jelikož neexistuje)
% U dizertační práce jsou typicky dva až tři oponenti. Pokud je chcete mít na titulním slajdu, prosím ručně odkomentujte a upravte jejich jména v definici "VUT title page" v souboru thesis.sty.
\opponent[Ing.]{Lukáš}{Pohl}[Ph.D.]

%%% Název práce
%  Parametr ve složených závorkách {} je název v originálním jazyce,
%  parametr v hranatých závorkách [] je překlad (podle toho jaký je originální jazyk).
%  V případě, že název Vaší práce je dlouhý a nevleze se celý do zápatí prezentace, použijte příkaz
%  \def\insertshorttitle{Zkác.\ náz.\ práce}
%  kde jako parametr vyplníte zkrácený název. Pokud nechcete zkracovat název, budete muset předefinovat,
%  jak se vytváří patička slidu. Viz odkaz: https://bit.ly/3EJTp5A
\title[Field-oriented control of five-phase permanent magnet synchronous motor]{Vektorové řízení pětifázového synchronního motoru s permanentními magnety}

%%% Označení oboru studia
%  Parametr ve složených závorkách {} je název oboru v originálním jazyce,
%  parametr v hranatých závorkách [] je překlad
\specialization[Teleinformatics]{Teleinformatika}

%%% Označení ústavu
%  Parametr ve složených závorkách {} je název ústavu v originálním jazyce,
%  parametr v hranatých závorkách [] je překlad
\department[Department of Control and Instrumentation]{Ústav automatizace a měřicí techniky}
%\department[Department of Biomedical Engineering]{Ústav biomedicínského inženýrství}
%\department[Department of Electrical Power Engineering]{Ústav elektroenergetiky}
%\department[Department of Electrical and Electronic Technology]{Ústav elektrotechnologie}
%\department[Department of Physics]{Ústav fyziky}
%\department[Department of Foreign Languages]{Ústav jazyků}
%\department[Department of Mathematics]{Ústav matematiky}
%\department[Department of Microelectronics]{Ústav mikroelektroniky}
%\department[Department of Radio Electronics]{Ústav radioelektroniky}
%\department[Department of Theoretical and Experimental Electrical Engineering]{Ústav teoretické a experimentální elektrotechniky}
%\department[Department of Telecommunications]{Ústav telekomunikací}
%\department[Department of Power Electrical and Electronic Engineering]{Ústav výkonové elektrotechniky a elektroniky}

%%% Označení fakulty
%  Parametr ve složených závorkách {} je název fakulty v originálním jazyce,
%  parametr v hranatých závorkách [] je překlad
%\faculty[Faculty of Architecture]{Fakulta architektury}
\faculty[Faculty of Electrical Engineering and~Communication]{Fakulta elektrotechniky a~komunikačních technologií}
%\faculty[Faculty of Chemistry]{Fakulta chemická}
%\faculty[Faculty of Information Technology]{Fakulta informačních technologií}
%\faculty[Faculty of Business and Management]{Fakulta podnikatelská}
%\faculty[Faculty of Civil Engineering]{Fakulta stavební}
%\faculty[Faculty of Mechanical Engineering]{Fakulta strojního inženýrství}
%\faculty[Faculty of Fine Arts]{Fakulta výtvarných umění}
%
%Nastavení logotypu (v hranatych zavorkach zkracene logo, ve slozenych plne):
\facultylogo[logo/FEKT_zkratka_barevne_PANTONE_CZ]{logo/png2pdf}

%%% Rok odevzdání práce
\graduateyear{2026}
%%% Akademický rok odevzdání práce
\academicyear{2025/26}

%%% Datum obhajoby (uplatní se pouze v prezentaci k obhajobě)
\date{10.\,06.\,2026} 

%%% Místo obhajoby
% Na titulních stránkách bude automaticky vysázeno VELKÝMI písmeny (pokud tyto stránky sází šablona)
\city{Brno}

%%% Abstrakt
\abstract[%
This thesis presents the design and implementation of field-oriented control of a five-phase interior permanent magnet synchronous motor (IPMSM).

Five-phase motors achieve higher torque density, smoother torque ripple, and greater fault tolerance compared to their three-phase equivalents. To describe these motors, an extended Clarke and Park transformation into the $dq_{1\,3}$ space is introduced, decomposing motor quantities into two parallel planes --- $dq_1$ capturing the fundamental harmonic component and $dq_3$ capturing its third harmonic. A motor model suitable for controller design is derived using this transformation.

The motor is driven by space vector pulse width modulation (SVPWM) with a symmetric vector arrangement. A cascade control structure is designed, consisting of an inner current loop and an outer speed loop. Back-EMF compensation is verified in the current loop, reducing overshoots during step changes of the reference value.

Reference currents are generated by the Maximum Torque Per Ampere (MTPA) algorithm, maximising the torque produced per ampere drawn from the supply. Third harmonic injection is also implemented, increasing the theoretical maximum torque by approximately 20\,\% compared to control utilising only the fundamental harmonic subspace. Current references are precomputed into a lookup table to reduce the online computational burden.

The designed algorithms are implemented and verified in Matlab/Simulink, both on a custom analytical model and on a Simscape model. The control blocks are subsequently adapted for HDL code generation using the Matlab HDL Coder and HDL Verifier toolboxes. The generated IP cores are deployed and functionally verified on the ZedBoard development kit (SoC Xilinx Zynq-7000). The implemented solution utilises approximately 90\,\% of the available FPGA logic resources.
]{%
Diplomová práce se zabývá návrhem a implementací vektorového řízení pětifázového
synchronního motoru s permanentními magnety umístěnými uvnitř rotoru (IPMSM).

Pětifázové motory oproti třífázovým dosahují vyšší hustoty točivého momentu,
plynulejšího průběhu a větší odolnosti vůči poruchám. Pro jejich popis je
představena rozšířená Clarkova a Parkova transformace do prostoru $dq_{1\,3}$,
která motorové veličiny rozděluje do dvou rovnoběžných rovin -- $dq_1$ zachycující
1. harmonickou složku a $dq_3$ zachycující její trojnásobek. Pomocí této
transformace je odvozen model motoru vhodný pro návrh regulátorů.

Pro buzení motoru je použita prostorově vektorová pulzně šířková modulace (SVPWM)
se symetrickým uspořádáním vektorů. Je navrženo kaskádní řízení složené
z vnitřní proudové smyčky a vnější smyčky otáčkové. V proudové smyčce je ověřena
kompenzace zpětně indukovaného elektromotorického napětí, která snižuje překmity
při skokové změně žádané hodnoty.

Referenční proudy jsou generovány algoritmem MTPA (Maximum Torque Per Ampere),
zajišťujícím maximální moment na odebíraný ampér ze zdroje. Dále je implementována
injekce 3. harmonické, která teoretický maximální moment zvyšuje přibližně
o 20\,\% oproti řízení využívajícímu pouze prostor 1. harmonické. Proudové
reference jsou předpočítány do vyhledávací tabulky, čímž se snižuje výpočetní
náročnost za běhu.

Navržené algoritmy jsou implementovány a ověřeny v prostředí Matlab/Simulink,
a to jak na vlastním analytickém modelu, tak na modelu ze knihovny Simscape.
Řídicí bloky jsou následně přizpůsobeny pro generování HDL kódu pomocí Matlab
HDL Coder a HDL Verifier toolboxů. Vygenerovaná IP jádra jsou nahrána
a funkčně ověřena na vývojovém kitu ZedBoard (SoC Xilinx Zynq-7000).
Navržené řešení využívá přibližně 90\,\% dostupných logických zdrojů FPGA.

  }

%%% Klíčová slova
\keywrds[ 
  field oriented controll, spacevector pulse width modulation, 5pahse permanent magnet synchronous motor
]{%
vektorové řízení, 5-tifázová motor s permantními magnety, prostorově voektorová pulzně šířková modulce
}

%%% Poděkování
\acknowledgement{%
Rád bych poděkoval vedoucímu diplomové práce
panu Ing.~Liboru Veslému, Ph.D.\ za odborné vedení,
konzultace, trpělivost a~podnětné návrhy k~práci.
}%